\uuid{HgaG}
\exo7id{3662}
\titre{exo7 3662}
\auteur{quercia}
\organisation{exo7}
\datecreate{2010-03-11}
\isIndication{false}
\isCorrection{true}
\chapitre{Espace euclidien, espace normé}
\sousChapitre{Produit scalaire, norme}
\cpge{mpsi-mp2i.espaces-prehilbertiens.a, pcsi.espaces-prehilbertiens.a, ptsi.geometrie-plan.b, mpsi-mp2i.espaces-prehilbertiens.b, pcsi.espaces-prehilbertiens.b, ptsi.geometrie-espace.b, mpsi-mp2i.espaces-prehilbertiens.c, pcsi.espaces-prehilbertiens.c, mpsi-mp2i.espaces-prehilbertiens.d, pcsi.espaces-prehilbertiens.d}
\module{Algèbre}
\niveau{L2}
\difficulte{}

\contenu{
\question{
Soit $E  =  \R_2[X]$ muni du produit scalaire :
$(P\mid Q) = \sum_{i=0}^4 P(i)Q(i)$.
Chercher une base orthonormée de $E$.
}
\reponse{
$\frac 1{\sqrt5}, \frac {X-2}{\sqrt{10}}, \frac {X^2-4X+2}{\sqrt{14}}$.
}
}
